\documentclass[10pt,a4paper]{amsart}
\usepackage[margin=19mm]{geometry}
\usepackage{amsmath,amssymb,mathtools}
\usepackage[scr=rsfs,cal=euler]{mathalfa}
\usepackage{xfrac}
\usepackage[unicode,hidelinks,bookmarks=false]{hyperref}
\newcommand{\ie}{i.e.\ }
\newcommand{\cf}{cf.\ }
\newcommand{\resp}{respectively\ }
\newcommand{\Subsection}{\subsection}
\providecommand{\nobreakdash}{\nobreak\hbox{-}}
\setlength{\emergencystretch}{2em}
\multlinegap0pt
\usepackage{xcolor}
\usepackage{mathtools}
\usepackage{amssymb}
\usepackage[shortlabels]{enumitem}
\usepackage{tikz}
\newtheorem{prop}{Proposition}
\DeclareMathOperator{\PW}{PW}
\title{Helson Inequality, Hankel Operators, and Weak Factorization on Paley-Wiener spaces of Convex Domains}
\author{Konstantinos Bampouras}
\date{Source: arXiv:2609.10730v1 (CC BY 4.0)}
\begin{document}
\maketitle

\noindent\emph{Source and licence.} Helson Inequality, Hankel Operators, and Weak Factorization on Paley-Wiener spaces of Convex Domains. Version arXiv:2609.10730v1 (CC BY 4.0), distributed by arXiv under the Creative Commons Attribution 4.0 International licence, http://creativecommons.org/licenses/by/4.0/, as recorded in the arXiv licence metadata for this exact version and shown on the abstract page https://arxiv.org/abs/2609.10730v1. Full licence notice: \texttt{licenses/arxiv-2609.10730-CC-BY-4.0.txt}.\par\smallskip

\noindent\textit{Editorial scope.} Counted unit: the article's prop completion (source0.tex lines 470--476). The article's complete original proof of it is delivered with it (source0.tex lines 477--488, one \texttt{proof} environment of 12 lines). No mathematics is written, repaired or re-derived: the statement and the proof are the article's own lines, in the article's own notation, and no retained line is substituted or re-flowed.  No further source-local context is needed: every cross-reference of every retained line resolves inside the two delivered spans.  Nothing was omitted from any retained span, so the package carries no omission record.  No retained line contains a citation, so the wrapper carries no bibliography and no bibliography entry.  No retained line contains a figure environment, an image-inclusion command or drawing code, so no image is delivered and the package declares no asset dependency.  Editorial formatting only, in addition: an AMS \texttt{amsart} preamble that restates the article's own declarations and macros used by the retained lines, namely the environments \texttt{prop} on separate counters, together with the standard packages those lines need.  The retained environments print in this document as Proposition 1 in order of appearance, whereas the article prints the same items with the numbering of its own full document.  The complete native source of the article as distributed in the arXiv e-print for this exact version is bundled unchanged as \texttt{sources/209944/source0.tex}.
\begin{prop}\label{completion}
	Let $\Omega$ be a convex set that does not contain affine lines and let $\mathcal{W}_p(\Omega)$, $1\leq p\leq \infty$, be the completion of $W_p(\Omega)$. Then, if $p<\infty$, $f\in \mathcal{W}_p(\Omega)$ if and only if there are orthogonal sequences $f_j,g_j\in \PW(\Omega)$ such that
	$$f=\lim_{N\to \infty} \sum_{j=1}^N f_j g_j \text{ in $\mathcal{W}_p(\Omega)$,} \quad \| (\|f_j\|_{L^2}\|g_j\|_{L^2})_{j}\|_{\ell^p}<\infty.$$
	Furthermore
	$$\|f\|_{\mathcal{W}_p(\Omega)}=\inf \|(\|f_j\|_{L^2}\|g_j\|_{L^2})_j\|_{\ell^p},$$ where the infimum is taken over all representations of $f$.
	For $p=\infty$ the proposition holds if we replace the condition $\| (\|f_j\|_{L^2}\|g_j\|_{L^2})_{j}\|_{\ell^\infty}<\infty$ by $\|f_j\|_{L^2}\|g_j\|_{L^2}\to 0$.
\end{prop}

\begin{proof}
	We will only treat the case $1\leq p<\infty.$ The case $p=\infty$ follows similarly, replacing $S^\infty$ by compact operators. For $f\in W_p(\Omega)$, let us define for a representation of $f=\sum_j f_jg_j$ the conjugate linear operator $T:\PW(\Omega)\to \PW(\Omega)$ as
	$$Tx=\sum_j \langle f_j,x\rangle g_j.$$
Let also $e_1,e_2,...$ be an orthonormal basis for $\PW(\Omega)$. Then, we can notice that
	$$\|f\|_{W_p(\Omega)}=\inf\{\|T\|_{S^p}: T\text{ is finite rank conjugate linear with }\sum_j T(e_j)e_j=f\}.$$
	It is therefore evident that
	$$\|f\|_{\mathcal{W}_p(\Omega)}=\inf\{\|T\|_{S^p}: T\in S^p\text{ conjugate linear with }\sum_j T(e_j)e_j=f \text{ in }\mathcal{W}_p(\Omega)\}.$$
	Since any such operator $T\in S^p$ can be written as 
	$$Tx=\sum_j \langle f_j,x\rangle g_j,$$ for $f_j,g_j\in \PW(\Omega)\setminus\{0\}$ orthogonal sequences, we have that
	$$\|T\|_{S^p}=\left(\sum_j (\|f_j\|_{L^2}\|g_j\|_{L^2})^p\right)^{1/p}.$$
	The proof is complete.
\end{proof}

\end{document}
